\section*{Chapter \ref{ch:b2:multint} --- Line Integrals and Multiple Integrals}

\begin{solution}{exo:b2:multint:1}
(a) Segment $\gamma(t) = (t, t)$, $t \in [0,1]$:
\[
\int_\gamma y^2\dd x + x\dd y
= \int_0^1 (t^2 + t)\,\dd t = \frac13 + \frac12 = \frac56 .
\]
(b) Parabola $\gamma(t) = (t, t^2)$:
\[
\int_0^1 \bigl(t^4\cdot 1 + t\cdot 2t\bigr)\dd t
= \frac15 + \frac23 = \frac{13}{15} .
\]
The two values differ, so the integral is path-dependent: the form
is \emph{not} exact --- consistently, $P_y = 2y \neq 1 = Q_x$, so
it is not even closed.
\end{solution}

\begin{solution}{exo:b2:multint:2}
$P = 2xy + y^3$, $Q = x^2 + 3xy^2 + 1$: $P_y = 2x + 3y^2 = Q_x$,
closed on $\R^2$. Search for $f$ with $f_x = P$: $f = x^2y + xy^3
+ g(y)$; then $f_y = x^2 + 3xy^2 + g'(y) = Q$ forces $g'(y) = 1$,
say $g(y) = y$. So
\[
f(x, y) = x^2y + xy^3 + y
\]
is a potential ($\R^2$ is star-shaped, so a potential had to
exist by Poincaré's lemma --- but exhibiting it is quicker). By
\cref{thm:b2:multint:ftc}, for any arc from $(0,0)$ to $(1,2)$:
\[
\int_\gamma\omega = f(1, 2) - f(0, 0) = 2 + 8 + 2 = 12 .
\]
\end{solution}

\begin{solution}{exo:b2:multint:3}
The domain is $\{0 \leq x \leq 1,\ x^2 \leq y \leq x\}$.
$y$-first:
\[
\int_0^1\!\int_{x^2}^{x}(x + y)\,\dd y\,\dd x
= \int_0^1\Bigl(x(x - x^2)
+ \frac{x^2 - x^4}{2}\Bigr)\dd x
= \int_0^1\Bigl(\frac{3x^2}{2} - x^3 - \frac{x^4}{2}\Bigr)\dd x
= \frac12 - \frac14 - \frac1{10} = \frac{3}{20} .
\]
$x$-first: the slice at height $y \in [0, 1]$ is $y \leq x \leq
\sqrt y$, so
\[
\int_0^1\!\int_{y}^{\sqrt y}(x + y)\,\dd x\,\dd y
= \int_0^1\Bigl(\frac{y - y^2}{2}
+ y(\sqrt y - y)\Bigr)\dd y
= \frac14 - \frac16 + \frac25 - \frac13 = \frac{3}{20} .
\]
\end{solution}

\begin{solution}{exo:b2:multint:4}
\emph{First integral.} On the disk $D_R$, in polar coordinates:
\[
\iint_{D_R}\frac{\dd x\,\dd y}{(1 + x^2 + y^2)^2}
= \int_0^{2\pi}\!\!\int_0^R
\frac{\rho\,\dd\rho\,\dd\alpha}{(1 + \rho^2)^2}
= 2\pi\Bigl[-\frac{1}{2(1 + \rho^2)}\Bigr]_0^R
= \pi\Bigl(1 - \frac{1}{1 + R^2}\Bigr)
\xrightarrow[R \to \infty]{} \pi .
\]
\emph{Second integral.} The quarter disk is $0 \leq \alpha \leq
\frac\pi2$, $0 \leq \rho \leq 1$, and $xy =
\rho^2\cos\alpha\sin\alpha$:
\[
\iint_{D'}xy\,\dd x\,\dd y
= \int_0^{\pi/2}\!\!\cos\alpha\sin\alpha\,\dd\alpha
\int_0^1 \rho^3\,\dd\rho
= \frac12\cdot\frac14 = \frac18 .
\]
\end{solution}

\begin{solution}{exo:b2:multint:5}
By \cref{cor:b2:multint:area} with $x = \cos^3 t$, $y = \sin^3 t$:
$x' = -3\cos^2 t\sin t$, $y' = 3\sin^2 t\cos t$, so
\[
xy' - yx' = 3\cos^4 t\sin^2 t + 3\sin^4 t\cos^2 t
= 3\sin^2 t\cos^2 t = \frac{3}{4}\sin^2 2t
= \frac{3}{8}(1 - \cos 4t) .
\]
Hence
\[
A = \frac12\int_0^{2\pi}\frac38(1 - \cos 4t)\,\dd t
= \frac{3}{16}\cdot 2\pi = \frac{3\pi}{8} .
\]
(The astroid fits in the unit disk of area $\pi$; three eighths of
$\pi$ is plausible for its four-cusped star shape.)
\end{solution}

\begin{solution}{exo:b2:multint:6}
\emph{Stacking:} above each $(x, y)$ of the unit disk $D$, $z$
runs from $x^2 + y^2$ to $1$:
\[
V = \iint_D \bigl(1 - x^2 - y^2\bigr)\dd x\,\dd y
= \int_0^{2\pi}\!\!\int_0^1 (1 - \rho^2)\rho\,\dd\rho\,\dd\alpha
= 2\pi\Bigl(\frac12 - \frac14\Bigr) = \frac\pi2 .
\]
\emph{Slicing:} the slice at height $z \in [0, 1]$ is the disk
$x^2 + y^2 \leq z$, of area $\pi z$:
\[
V = \int_0^1 \pi z\,\dd z = \frac\pi2 .
\]
\end{solution}

\begin{solution}{exo:b2:multint:7}
In spherical coordinates the volume element is
$r^2\cos\varphi\,\dd r\,\dd\theta\,\dd\varphi$
(\cref{ex:b2:multint:spherical}), and the angular part integrates
to $4\pi$ ($2\pi$ from $\theta$, $\int_{-\pi/2}^{\pi/2}\cos = 2$).
The shell volume is
\[
\int_a^b 4\pi r^2\,\dd r = \frac43\pi\bigl(b^3 - a^3\bigr).
\]
For the second integral, the integrand $1/r$ depends only on $r$:
\[
\iiint_B \frac{\dd x\,\dd y\,\dd z}{r}
= \int_0^R 4\pi r^2\cdot\frac1r\,\dd r
= 4\pi\,\frac{R^2}{2} = 2\pi R^2 .
\]
(The integrand blows up at the origin, but harmlessly: $r^2/r = r$
is continuous --- the integral over the shells $\varepsilon \leq r
\leq R$ converges as $\varepsilon \to 0$, which is the precise
sense of the statement. This kind of computation is the first
step toward Newton's theorem that a homogeneous ball attracts like
a point mass at its center.)
\end{solution}

\begin{solution}{exo:b2:multint:8}
The integrand $g(t; x, y) = xP(tx, ty) + yQ(tx, ty)$ is
$\mathcal{C}^1$ in $(x, y)$, continuous in $t$, with partial
derivatives continuous on $[0,1] \times U$; differentiation under
the integral sign (\cref{ch:b2:integration}, applied on the
compact $t$-interval $[0,1]$, domination being automatic there)
gives
\[
f_x(x, y)
= \int_0^1 \bigl(P(tx, ty) + tx\,P_x(tx, ty)
+ ty\,Q_x(tx, ty)\bigr)\dd t .
\]
Using closedness $Q_x = P_y$:
\[
tx\,P_x(tx, ty) + ty\,P_y(tx, ty)
= t\,\frac{\dd}{\dd t}\bigl[P(tx, ty)\bigr] ,
\]
so the integrand is $P(tx, ty) + t\frac{\dd}{\dd t}P(tx, ty) =
\frac{\dd}{\dd t}\bigl[t\,P(tx, ty)\bigr]$ and
\[
f_x(x, y) = \Bigl[t\,P(tx, ty)\Bigr]_0^1 = P(x, y) .
\]
Symmetrically $f_y = Q$ (same computation with $P_y = Q_x$ used
the other way). Note where the hypothesis enters: $f$ is defined
by integrating along the segment $[0, M]$, which lies in $U$
precisely because $U$ is star-shaped.
\end{solution}

\begin{solution}{exo:b2:multint:9}
On the strip $[0, A] \times [0, \infty)$ the function $(x, y)
\mapsto e^{-xy}\sin x$ is not absolutely integrable up to $y =
\infty$ uniformly in a naive sense, but each iterated integral
converges and their equality follows from Fubini on $[0, A]
\times [0, B]$ plus a limit $B \to \infty$ (the tail
$\int_0^A\int_B^\infty e^{-xy}\abs{\sin x}\,\dd y\,\dd x \leq
\int_0^A \frac{e^{-Bx}\abs{\sin x}}{x}\dd x \leq \int_0^A e^{-Bx}
\dd x\to 0$, using $\abs{\sin x} \leq x$).

\emph{$y$ first:} $\int_0^\infty e^{-xy}\,\dd y = \frac1x$ for $x
> 0$, so the first integral is $\int_0^A \frac{\sin x}{x}\,\dd x$.

\emph{$x$ first:} two integrations by parts (or taking the
imaginary part of $\int_0^A e^{(i - y)x}\dd x$) give
\[
\int_0^A e^{-xy}\sin x\,\dd x
= \frac{1 - e^{-Ay}(y\sin A + \cos A)}{1 + y^2} .
\]
Integrating in $y$ over $[0, \infty)$, the term $\int_0^\infty
\frac{\dd y}{1 + y^2} = \frac\pi2$ splits off:
\[
\int_0^A \frac{\sin x}{x}\,\dd x
= \frac{\pi}{2}
- \int_0^\infty e^{-Ay}\,\frac{y\sin A + \cos A}{1 + y^2}\,\dd y .
\]
The remainder is bounded by $\int_0^\infty e^{-Ay}\frac{y +
1}{1 + y^2}\dd y \leq \int_0^\infty e^{-Ay}\cdot\frac{1+y}{1+y^2}
\,\dd y \to 0$ as $A \to \infty$ (dominated convergence, or the
crude bound $\frac{1 + y}{1 + y^2} \leq \frac32$ giving
$\frac{3}{2A}$). Hence $\int_0^\infty\frac{\sin x}{x}\dd x =
\frac\pi2$ --- the same value obtained in
\cref{ch:b2:integration} by differentiating a parameter integral;
here Fubini does the work instead.
\end{solution}

\begin{solution}{exo:b2:multint:10}
Parametrize by arc length $s \in [0, 2\pi]$, so $x'^2 + y'^2 = 1$,
and translate so that $\int_0^{2\pi} x(s)\,\dd s = 0$. By
\cref{cor:b2:multint:area},
\[
2A = \oint x\,\dd y - y\,\dd x
= \int_0^{2\pi}\bigl(xy' - yx'\bigr)\dd s .
\]
Integrating $\oint y\,\dd x$ by parts over the period (boundary
terms cancel by periodicity), $-\int yx' = \int y'x$, so in fact
$2A = 2\int_0^{2\pi}xy'\,\dd s$. Then $2xy' \leq x^2 + y'^2$
gives
\[
2A \leq \int_0^{2\pi}\bigl(x^2 + y'^2\bigr)\dd s
= \int_0^{2\pi} x^2 + \int_0^{2\pi}\bigl(1 - x'^2\bigr)
= 2\pi - \int_0^{2\pi}\bigl(x'^2 - x^2\bigr)\dd s .
\]
Wirtinger's inequality (\cref{ch:b2:fourier} exercises: for a
$2\pi$-periodic $\mathcal{C}^1$ function with zero mean,
$\int x^2 \leq \int x'^2$) makes the last integral nonnegative:
$A \leq \pi$. Equality requires equality in Wirtinger ($x(s) =
a\cos s + b\sin s$) and in $2xy' \leq x^2 + y'^2$ ($y' = x$
pointwise), which forces $y = a\sin s - b\cos s + c$: the curve is
the unit circle (suitably centered). Since a curve of length $L$
rescales to length $2\pi$, the general statement is $A \leq
\frac{L^2}{4\pi}$: among all closed curves of given perimeter, the
circle encloses the largest area.
\end{solution}

\begin{solution}{exo:b2:multint:11}
In spherical coordinates $z = r\sin\varphi$ and $\dd x\,\dd
y\,\dd z = r^2\cos\varphi\,\dd r\,\dd\theta\,\dd\varphi$:
\[
\iiint_B z^2
= \int_0^R r^4\,\dd r\int_0^{2\pi}\dd\theta
\int_{-\pi/2}^{\pi/2}\sin^2\varphi\cos\varphi\,\dd\varphi
= \frac{R^5}5\cdot2\pi\cdot
\Bigl[\frac{\sin^3\varphi}3\Bigr]_{-\pi/2}^{\pi/2}
= \frac{4\pi R^5}{15}.
\]
By the symmetry of the ball under permuting coordinates,
$\iiint_B x^2 = \iiint_B y^2 = \iiint_B z^2$, so $\iiint_B(x^2
+ y^2 + z^2) = 3\cdot\frac{4\pi R^5}{15} = \frac{4\pi R^5}5$.
Shell check: $\int_0^R r^2\cdot 4\pi r^2\,\dd r = \frac{4\pi
R^5}5$ --- the integrand $r^2$ is constant on the sphere of
radius $r$, of area $4\pi r^2$.
\end{solution}

\begin{solution}{exo:b2:multint:12}
Parametrize the edge from $(x_i, y_i)$ to $(x_{i+1}, y_{i+1})$
by $\gamma(t) = \bigl((1-t)x_i + tx_{i+1},\ (1-t)y_i +
ty_{i+1}\bigr)$. Its contribution to
$\frac12\oint(x\,\dd y - y\,\dd x)$ is
\[
\frac12\int_0^1\Bigl(\bigl((1-t)x_i +
tx_{i+1}\bigr)(y_{i+1} - y_i) - \bigl((1-t)y_i +
ty_{i+1}\bigr)(x_{i+1} - x_i)\Bigr)\dd t ,
\]
and since $\int_0^1\bigl((1-t)u + tv\bigr)\dd t = \frac{u +
v}2$, this equals
\[
\frac14\Bigl((x_i + x_{i+1})(y_{i+1} - y_i) - (y_i +
y_{i+1})(x_{i+1} - x_i)\Bigr)
= \frac12\bigl(x_iy_{i+1} - x_{i+1}y_i\bigr),
\]
the cross terms cancelling. Summing over the $m$ edges gives
the shoelace formula, by \cref{cor:b2:multint:area}. Triangle
$(0,0), (1,0), (0,1)$: $\frac12\bigl((0\cdot0 - 1\cdot0) +
(1\cdot1 - 0\cdot0) + (0\cdot0 - 0\cdot1)\bigr) = \frac12$,
the correct area.
\end{solution}

\begin{solution}{pb:b2:multint:1}
\textbf{1.} The ball $B_n(R)$ is described by iterated bounds
$-R \leq x_n \leq R$, then $\abs{x_{n-1}} \leq \sqrt{R^2 -
x_n^2}$, and so on; substituting $x_i = Ru_i$ in each of the
$n$ one-variable integrals multiplies each by $R$ and maps
the bounds onto those of $B_n(1)$: $v_n(R) = R^n\,v_n(1) =
V_nR^n$.

\textbf{2.} Slicing along $x_n = t$: the slice of $B_n(1)$ is
the ball $B_{n-1}\bigl(\sqrt{1 - t^2}\bigr)$, so, by question
1,
\[
V_n = \int_{-1}^1 v_{n-1}\bigl(\sqrt{1 - t^2}\bigr)\,\dd t
= V_{n-1}\int_{-1}^{1}(1 - t^2)^{\frac{n-1}2}\,\dd t .
\]

\textbf{3.} With $t = \sin\theta$, $\dd t =
\cos\theta\,\dd\theta$ and $(1 - t^2)^{\frac{n-1}2} =
\cos^{n-1}\theta$ on $\intcc{-\pi/2}{\pi/2}$:
\[
\int_{-1}^1(1 - t^2)^{\frac{n-1}2}\dd t
= \int_{-\pi/2}^{\pi/2}\cos^n\theta\,\dd\theta
= 2\int_0^{\pi/2}\cos^n\theta\,\dd\theta = 2W_n,
\]
and $\theta \mapsto \frac\pi2 - \theta$ interchanges the sine
and cosine forms of $W_n$.

\textbf{4.} Write $\sin^n = \sin^{n-2}(1 - \cos^2)$ and
integrate $\int\sin^{n-2}\cos\cdot\cos$ by parts ($v =
\frac{\sin^{n-1}}{n-1}$):
\[
W_n = W_{n-2} - \frac{W_n}{n-1}
\quad\Longrightarrow\quad
W_n = \frac{n-1}{n}W_{n-2}.
\]
Hence $nW_nW_{n-1} = (n-1)W_{n-1}W_{n-2}$: the sequence
$(nW_nW_{n-1})$ is constant, equal to $1\cdot W_1W_0 =
1\cdot\frac\pi2$, so $W_nW_{n-1} = \frac{\pi}{2n}$.

\textbf{5.} Questions 2--3 give $V_n = 2W_nV_{n-1}$, twice:
\[
V_n = 2W_n\cdot 2W_{n-1}\,V_{n-2}
= 4\,\frac{\pi}{2n}\,V_{n-2} = \frac{2\pi}n\,V_{n-2}.
\]

\textbf{6.} From $V_0 = 1$: $V_{2k} = \frac{2\pi}{2k}V_{2k-2}
= \frac\pi kV_{2k-2}$, so $V_{2k} = \frac{\pi^k}{k!}$ by
induction. From $V_1 = 2$: $V_{2k+1} =
\frac{2\pi}{2k+1}V_{2k-1}$, so
\[
V_{2k+1} = 2\prod_{j=1}^k\frac{2\pi}{2j+1}
= \frac{2^{k+1}\pi^k}{1\cdot3\cdot5\cdots(2k+1)} .
\]

\textbf{7.} $V_1 = 2$, $V_2 = \pi \approx 3.142$, $V_3 =
\frac{4\pi}3 \approx 4.189$, $V_4 = \frac{\pi^2}2 \approx
4.935$, $V_5 = \frac{8\pi^2}{15} \approx 5.264$, $V_6 =
\frac{\pi^3}6 \approx 5.168$, $V_7 = \frac{16\pi^3}{105}
\approx 4.725$. The one-step ratio is $V_n/V_{n-1} = 2W_n$,
and $(W_n)$ is decreasing ($\sin^n \leq \sin^{n-1}$
pointwise). Now $2W_5 = 2\cdot\frac45\cdot\frac23 =
\frac{16}{15} > 1$ while $2W_6 =
2\cdot\frac56\cdot\frac34\cdot\frac12\cdot\frac\pi2 =
\frac{5\pi}{16} < 1$: the ratios exceed $1$ up to $n = 5$ and
are below $1$ from $n = 6$ on --- $(V_n)$ increases to its
maximum $V_5$ and then decreases.

\textbf{8.} For $n \geq 13 > 4\pi$: $V_n/V_{n-2} = 2\pi/n <
\tfrac12$, so $V_{n} \leq C\cdot 2^{-n/2}$ with a fixed
constant; better, for any $q > 0$, $2\pi/n < q^2$ for large
$n$, so $V_n/q^n \to 0$: the decay beats every geometric
sequence. Convergence of $\sum V_n$ follows from the ratio
$V_n/V_{n-2} \to 0$ (compare with a geometric series from
some rank on).

\textbf{9.} $\sum_{k\geq0}V_{2k}x^{2k} =
\sum_{k\geq0}\frac{(\pi x^2)^k}{k!} = \eu^{\pi x^2}$, the
exponential series (\cref{ch:b2:powerseries}), convergent for
every $x$. At $x = 1$: $\sum_kV_{2k} = \eu^\pi \approx
23.14$.

\textbf{10.} On the cube $\intcc{-R}R^n$ the integrand is the
product $\prod_i\eu^{-x_i^2}$, so the iterated integral
factorizes: $\bigl(\int_{-R}^R\eu^{-t^2}\dd t\bigr)^n$.
Letting $R \to \infty$ and using
$\int_\R\eu^{-t^2}\dd t = \sqrt\pi$
(\cref{ex:b2:multint:polar}): $I_n = \pi^{n/2}$.

\textbf{11.} $t = u^2$ gives $\Gamma(\tfrac12) =
\int_0^\infty t^{-1/2}\eu^{-t}\dd t =
2\int_0^\infty\eu^{-u^2}\dd u = \sqrt\pi$. Iterating
$\Gamma(s+1) = s\Gamma(s)$: $\Gamma(k+1) = k!\,\Gamma(1) =
k!$, and
\[
\Gamma\Bigl(k + \frac32\Bigr)
= \Bigl(k + \frac12\Bigr)\Bigl(k - \frac12\Bigr)\cdots
\frac12\cdot\Gamma\Bigl(\frac12\Bigr)
= \frac{(2k+1)(2k-1)\cdots1}{2^{k+1}}\,\sqrt\pi .
\]

\textbf{12.} Set $F_n = \pi^{n/2}/\Gamma(\frac n2 + 1)$.
Since $\Gamma(\frac n2 + 1) = \frac n2\,\Gamma(\frac n2) =
\frac n2\,\Gamma(\frac{n-2}2 + 1)$, we get $F_n =
\frac{2\pi}nF_{n-2}$: the same recursion as $V_n$ (question
5). Bases: $F_1 = \sqrt\pi/\Gamma(\frac32) =
\sqrt\pi/(\frac{\sqrt\pi}2) = 2 = V_1$ and $F_2 =
\pi/\Gamma(2) = \pi = V_2$. By induction $V_n = F_n$ for all
$n$; question 11 turns this back into the two closed forms of
question 6.

\textbf{13.} With $r = \sqrt t$: $\int_0^\infty
\eu^{-r^2}r^{n-1}\dd r = \frac12\int_0^\infty
t^{\frac n2 - 1}\eu^{-t}\dd t = \frac12\Gamma(\frac n2)$.
Hence
\[
n\,V_n\int_0^\infty\eu^{-r^2}r^{n-1}\dd r
= V_n\cdot\frac n2\,\Gamma\Bigl(\frac n2\Bigr)
= V_n\,\Gamma\Bigl(\frac n2 + 1\Bigr) = \pi^{n/2} = I_n .
\]
Both sides being proved, the identity can be \emph{read} as
the shell decomposition of the Gaussian integral: the sphere
of radius $r$ carries area $nV_nr^{n-1}$, and the Gaussian
weight $\eu^{-r^2}$ is integrated over the shells.

\textbf{14.} $s_0 = V_1 = 2$ (the $0$-sphere is two points),
$s_1 = 2V_2 = 2\pi$, $s_2 = 3V_3 = 4\pi$, $s_3 = 4V_4 =
2\pi^2$; and $\int_0^R s_{n-1}r^{n-1}\dd r = V_nR^n =
v_n(R)$: area is the radial derivative of volume.

\textbf{15.} For $n = 2k$, Stirling
(\cref{thm:b2:comparison:stirling}) gives $k! \sim
\sqrt{2\pi k}\,(k/\eu)^k$, so
\[
V_{2k} = \frac{\pi^k}{k!}
\sim \frac{(\pi\eu/k)^k}{\sqrt{2\pi k}}
= \frac1{\sqrt{\pi n}}\Bigl(\frac{2\pi\eu}n\Bigr)^{n/2}
\qquad (n = 2k).
\]
For odd $n$: $V_{2k+1} = 2W_{2k+1}V_{2k} \leq 2V_{2k}$, so
the same super-geometric decay bounds hold (up to a factor
$2$ and a shift of one in the exponent) --- for every $q >
0$, $V_n = o(q^n)$.

\textbf{16.} $V_2/4 = \pi/4 \approx 0.785$; $V_3/8 = \pi/6
\approx 0.524$; $V_{10}/2^{10} = \frac{\pi^5}{120\cdot1024}
\approx 0.0025$. In general $\frac{V_n/2^n}{V_{n-2}/2^{n-2}}
= \frac{2\pi}{4n} = \frac{\pi}{2n} \to 0$: the ratio tends to
$0$ (super-geometrically). The inscribed ball occupies a
vanishing fraction: the cube's volume migrates to its
corners.

\textbf{17.} By question 1 the inner ball of radius $1 -
\varepsilon$ has volume $V_n(1-\varepsilon)^n$, so the outer
shell carries the fraction $1 - (1 - \varepsilon)^n \to 1$.
For $n = 100$, $\varepsilon = 0.01$: $(0.99)^{100} =
\eu^{100\ln0.99} \approx \eu^{-1.005} \approx 0.366$: about
$63\%$ of the ball lies within $1\%$ of its surface.

\textbf{18.} $(W_n)$ decreases, so $W_n \leq W_{n-1} \leq
W_{n-2} = \frac{n}{n-1}W_n$: squeezing, $W_{n-1}/W_n \to 1$.
Multiplying by $W_nW_{n-1} = \frac\pi{2n}$: $W_n^2 \sim
\frac\pi{2n}$, i.e.\ $W_n \sim \sqrt{\pi/(2n)}$. Lower bound:
$W_n^2 \geq W_nW_{n+1} = \frac{\pi}{2(n+1)}$, so $W_n \geq
\sqrt{\pi/(2(n+1))}$ for every $n$.

\textbf{19.} Numerator: $1 - x^2 \leq \eu^{-x^2}$ gives $(1 -
x^2)^{\frac{n-1}2} \leq \eu^{-(n-1)x^2/2}$, and with $a =
\frac{n-1}2$,
\[
\int_\delta^1(1 - x^2)^{\frac{n-1}2}\dd x
\leq \int_\delta^\infty\eu^{-ax^2}\dd x
\leq \int_\delta^\infty\frac x\delta\,\eu^{-ax^2}\dd x
= \frac{\eu^{-a\delta^2}}{2a\delta}
= \frac{\eu^{-(n-1)\delta^2/2}}{(n-1)\delta}.
\]
Denominator: $2W_n \geq \sqrt{2\pi/(n+1)}$ by question 18.
Dividing gives the displayed bound. For $\delta =
s/\sqrt{n-1}$ it becomes
$\sqrt{\tfrac{n+1}{2\pi(n-1)}}\;\eu^{-s^2/2}/s = O\bigl(
\eu^{-s^2/2}/s\bigr)$, uniformly in $n$: outside the slab
$\abs{x_1} \leq s/\sqrt{n-1}$ there is almost no volume, for
$s$ moderately large --- and by symmetry the same holds for
every direction.

\textbf{20.} The two statements coexist because they describe
different coordinates of the same point. Almost every point
of $B_n(1)$ has norm close to $1$ (question 17: radial
concentration near the sphere), yet each of its $n$
coordinates is small, of order $1/\sqrt n$ (question 19),
which is consistent since $n$ coordinates of size $1/\sqrt n$
have norm of order $1$. High-dimensional volume concentrates
where all coordinates share the norm budget equally ---
near the sphere, but far from every coordinate axis pole.

\textbf{21.} Slice $\Delta_n$ at $x_n = t \in \intcc01$: the
slice is $\{x' \in \R^{n-1} : x_i \geq 0,\ \sum x_i \leq 1 -
t\} = (1-t)\Delta_{n-1}$, of volume
$(1-t)^{n-1}\operatorname{vol}(\Delta_{n-1})$ by homogeneity.
So
\[
\operatorname{vol}(\Delta_n) =
\operatorname{vol}(\Delta_{n-1})\int_0^1(1 - t)^{n-1}\dd t =
\frac{\operatorname{vol}(\Delta_{n-1})}{n}
\quad\Longrightarrow\quad
\operatorname{vol}(\Delta_n) = \frac1{n!}\,.
\]

\textbf{22.} The $2^n$ sign-orthants cut $C_n$ into $2^n$
copies of $\Delta_n$ (with negligible overlaps on coordinate
hyperplanes): $\operatorname{vol}(C_n) = \frac{2^n}{n!}$. If
$\sum\abs{x_i} \leq 1$ then $\sum x_i^2 \leq
\bigl(\sum\abs{x_i}\bigr)^2 \leq 1$: $C_n \subseteq B_n(1)$;
and $B_n(1) \subseteq \intcc{-1}1^n$ since $\abs{x_i} \leq
\norm x$. Hence $\frac{2^n}{n!} \leq V_n \leq 2^n$ ---
consistent with question 15, which places $V_n$ between the
factorial and geometric scales.

\textbf{23.} For $(x, y)$ in the unit disk, the slice of
$B_4(1)$ is the disk of radius $\sqrt{1 - x^2 - y^2}$ in the
$(z, w)$-plane, of area $\pi(1 - x^2 - y^2)$. In polar
coordinates:
\[
V_4 = \iint_{x^2+y^2\leq1}\pi(1 - x^2 - y^2)\,\dd x\,\dd y
= \pi\int_0^{2\pi}\!\!\int_0^1(1 - \rho^2)\rho\,
\dd\rho\,\dd\alpha
= \pi\cdot2\pi\cdot\frac14 = \frac{\pi^2}2 ,
\]
agreeing with question 6.

\textbf{24.} The probability is the volume ratio
$\dfrac{V_{20}}{2^{20}} = \dfrac{\pi^{10}}{10!\cdot2^{20}}
\approx \dfrac{0.0258}{1\,048\,576} \approx
2.5\cdot10^{-8}$. The number of draws until the first hit is
of order the inverse, about $4\cdot10^7$: a rejection sampler
that worked beautifully for the disk ($\pi/4$ of hits) is
useless in dimension $20$ --- the curse of dimensionality in
one line.

\textbf{25.} Route one (Parts I--II) used: Fubini-type
slicing of the iterated integral, one-variable substitution
in each coordinate (homogeneity), and the Wallis integrals
--- pure one-variable calculus plus induction. Route two
(Part III) used: Fubini for the product structure of $I_n$,
the polar change of variables through the Gauss integral of
\cref{ex:b2:multint:polar}, and the $\Gamma$ function's
functional equation. They meet in $V_n =
\pi^{n/2}/\Gamma(\frac n2 + 1)$, with Stirling
(\cref{thm:b2:comparison:stirling}) converting the formula
into asymptotics. The Year 3 volume rebuilds all of this on
Lebesgue's integral: there Fubini and the change of variables
are theorems for general integrable functions, spherical
coordinates exist in every dimension, and the same ball
volumes reappear as the worked dividends of the
product-measure and Stirling problems --- with dominated
convergence replacing our hand-made squeezes.
\end{solution}
